\section{Towards a Machine-Learned Interface (MLI) Predictor}

The ultimate goal of this pipeline is to develop a Machine-Learned Interface (MLI) predictor capable of estimating interface favourability directly from slab pair descriptors, without requiring full relaxation or enumeration of stacking shifts. This model will act as an early-stage filter to prioritise promising candidates for structure generation and DFT refinement.

Input features will include:
\begin{itemize}
    \item Surface energies
    \item Crystallographic orientation identifiers (Miller indices)
    \item Lattice mismatch metrics (based on in-plane scaling ratios)
    \item Interface registry scores (from ARTEMIS or RAFFLE shifts)
    \item Atomic coordination and stoichiometry at the intended interface plane
\end{itemize}

Initial labels for training will be derived from \( \Delta E_\mathrm{IF} \) values obtained via DFT or MACE (after benchmarking). Where available, these will be linked to structural features extracted from POSCAR files and relaxation trajectories.

Prediction targets will include:
\begin{itemize}
    \item Binary classification: favourable vs. unfavourable
    \item Regression: predicted interface energy \( \Delta E_\mathrm{IF} \)
    \item Rank scoring: predicted rank percentile within a batch
\end{itemize}

Simple models (e.g. decision trees, SVMs) will be used as baselines. If effective, graph-based architectures or equivariant neural networks may be explored to capture spatial features across slab boundaries.

The MLI predictor will integrate upstream with ARTEMIS and RAFFLE:
\begin{enumerate}
    \item Given a pair of slabs, the MLI will score interface likelihoods before full stack generation.
    \item High-scoring pairs will be passed to ARTEMIS for registry-matched stacking and to RAFFLE for disordered sampling.
    \item Final relaxed energies will serve to update or retrain the MLI in an active-learning loop.
\end{enumerate}

This workflow could enable a fully ML-driven interface discovery pipeline; one that uses minimal DFT data to bootstrap structure generation, ranking, and model refinement.

% TODO: Add schematic of MLI + ARTEMIS + RAFFLE integration pipeline

% ---------- %

\paragraph{MLI Model Architecture and Software Perspective}

The proposed Machine-Learned Interface (MLI) model will adopt a graph neural network (GNN) architecture, which is
particularly well-suited to atomic-scale prediction tasks.

In this formulation, atoms are represented as nodes, and edges encode interatomic relationships such as distance, bond
type, or spatial proximity.

Node features may include atomic number, electronegativity, or coordination number, while edge features capture pairwise
distances, bond angles, or periodic adjacency.

This graph-based abstraction mirrors the geometry and locality of real materials, enabling the model to learn chemically
and physically meaningful patterns across interface structures.

% ---------- %

GNNs are especially appropriate for this task because they respect the symmetries of atomistic systems; they are
invariant or equivariant to translation, rotation, and permutation of atoms, depending on their architecture.

Through message-passing and feature aggregation, each atom aggregates information from its neighbors to refine local
descriptors, allowing the model to learn both local bonding patterns and emergent long-range effects.

These properties are essential for predicting interfacial phenomena, which often arise from geometric asymmetry,
registry mismatch, and complex bonding environments.

% ---------- %

\paragraph{Implementation Pipeline and Integration}

From a software engineering standpoint, the MLI model will be implemented as a modular, callable component in the
interface discovery pipeline.

It will accept structured slab-pair data; either parsed directly from POSCARs or preprocessed descriptor arrays, and
produce predictions of interface favourability.

The model will be constructed using scalable graph learning libraries such as PyTorch Geometric, DGL, or JAX-based GNN
frameworks, and will provide a clean API for downstream tools such as ARTEMIS and RAFFLE. Outputs will be compatible
with JSON or HDF5-based metadata systems, ensuring traceability and reproducibility.

% ---------- %

Initial experiments may use relatively shallow GNNs with fixed neighbor cutoffs, but later iterations may incorporate
equivariant tensor representations (e.g. MACE-style angular embeddings) to better capture orientation-dependent effects.
Training targets will include both regression on \( \Delta E_\mathrm{IF} \) and binary classification
(favourable/unfavourable), optionally extended with uncertainty quantification.

Crucially, the MLI is not meant to replace DFT or full relaxation pipelines, but to act as a high-speed prefilter that
screen candidate structures for further evaluation.

If successful, it will form the decision engine in a closed-loop ML-driven workflow for interfacial discovery.